\setbeamertemplate{navigation symbols}{}
\setbeamercolor{background canvas}{bg=white}
\setbeamercolor{normal text}{fg=PInk}
\setbeamercolor{itemize item}{fg=PTeal}
\setbeamercolor{itemize subitem}{fg=PTeal}
\setbeamertemplate{itemize item}{\raisebox{0.28ex}{\rule{0.34em}{0.34em}}}
\setbeamertemplate{itemize subitem}{\raisebox{0.28ex}{\rule{0.26em}{0.26em}}}
\usebeamercolor[fg]{normal text}

% A block is a title, a rule and text. No fill, no border, no rounded corner,
% no shadow: hierarchy comes from the rule and from whitespace. That is the
% whole defence against looking like a 2012 conference poster.
\newenvironment{pblock}[1]{%
  \par\vspace{\blockskip}%
  {\color{PDeep}\bfseries\Large #1\par}%
  \vspace{2.5mm}{\color{PTeal}\rule{\linewidth}{1.5mm}}\par
  \vspace{6mm}\color{PInk}\justifying\ignorespaces
}{\par}

% The single filled variant, for at most one or two key results. Used more
% than that it stops being emphasis.
%
% Its own savebox rather than \@tempboxa: this file is \input in the preamble,
% where `@' is not a letter, so \@tempboxa would be read as \@ followed by the
% word "tempboxa" and fail with an obscure \spacefactor error.
\newsavebox{\pcalloutbox}
\newenvironment{pcallout}{%
  \par\vspace{\blockskip}%
  % \fboxsep is 4 mm globally for other boxes; left on, the colorbox adds it
  % outside the 8 mm padding below and the callout runs 8 mm past the column.
  \setlength{\fboxsep}{0pt}%
  \begin{lrbox}{\pcalloutbox}%
  \begin{minipage}{\dimexpr\linewidth-16mm\relax}\color{PInk}\justifying
}{%
  \end{minipage}\end{lrbox}%
  \colorbox{PMist}{\hspace{8mm}%
    \raisebox{0pt}[\dimexpr\ht\pcalloutbox+8mm\relax][\dimexpr\dp\pcalloutbox+8mm\relax]{%
      \usebox{\pcalloutbox}}\hspace{8mm}}\par}

% Figure caption. The `caption' package is incompatible with beamer and a
% poster has no floats, so this is all that is needed.
\newcommand{\figcap}[1]{\par\vspace{2.5mm}{\color{PSlate}\small #1\par}}

% Width is MANDATORY: the source PNGs carry inconsistent DPI metadata
% (100/140/150/160/170), so a bare \includegraphics sizes two similar figures
% differently. The \typeout line lets check.ps1 audit what was really applied.
\newcommand{\pfig}[2][\linewidth]{%
  \typeout{FIGCHECK: #2 width=\the\dimexpr#1\relax}%
  \includegraphics[width=#1]{#2}}

% Numbered steps and labelled definitions without beamer's enumerate or
% description environments, both of which recurse under enumitem.
\newcommand{\step}[2]{\par\vspace{0.4\baselineskip}%
  \noindent\hangindent=1.6em\hangafter=1%
  {\color{PTeal}\bfseries #1}\hspace{0.6em}#2\par}
\newcommand{\defn}[2]{\par\vspace{0.45\baselineskip}%
  \noindent{\color{PDeep}\bfseries #1}\par\nobreak\vspace{0.5mm}%
  \noindent #2\par}

\newcommand{\emphnum}[1]{{\color{PSand}\bfseries #1}}
\newcommand{\term}[1]{{\color{PDeep}\bfseries #1}}
\newcommand{\wip}{{\color{PSand}\bfseries\normalsize [\,in progress\,]}}
